\chapter{Συμπεράσματα}

Η παρούσα διπλωματική εργασία σχεδίασε, υλοποίησε και αξιολόγησε μια ολοκληρωμένη πλατφόρμα Ψηφιακού Διδύμου για σχολικές υποδομές, με πιλοτική εφαρμογή στα Εκπαιδευτήρια Πλάτων στην Κατερίνη. Το κεντρικό ερώτημα ήταν αν η ενοποίηση υλικού χαμηλού κόστους με λογισμικό ανοιχτού κώδικα μπορεί να δημιουργήσει έναν κλειστό βρόχο ελέγχου ικανό να βελτιώσει ταυτόχρονα το μαθησιακό περιβάλλον και τη λειτουργική διαχείριση ενός σχολικού κτιρίου. Τα ευρήματα παρουσιάζονται με ακαδημαϊκή ειλικρίνεια --- διαχωρίζοντας αυτό που αποδείχθηκε από αυτό που παρατηρήθηκε.

Το κύριο τεχνολογικό εύρημα είναι ότι η μετάβαση από παθητική παρακολούθηση (\textlatin{Digital Shadow}) σε αυθεντικό Ψηφιακό Δίδυμο είναι εφικτή με \textlatin{COTS} εξοπλισμό κόστους 245,17\euro. Η πλατφόρμα υλοποιεί τον πλήρη κύκλο \textlatin{Sense--Decide--Act} και περιλαμβάνει:

\begin{itemize}
\item Ακουστική παρακολούθηση με αυτόματη ενεργοποίηση \textlatin{quiet-lamp} βάσει ορίου 55 \textlatin{dBA}
\item Παρακολούθηση ψυκτικού εξοπλισμού με \textlatin{FridgeGuardian} και αυτόματες ειδοποιήσεις
\item \textlatin{Cloud SaaS} πλατφόρμα με \textlatin{3D visualization}, \textlatin{RBAC} και \textlatin{LLM-based energy forecasting}
\item \textlatin{Demo} περιβάλλον δημόσιας πρόσβασης στο \textlatin{digitaltwin.pi-tech.gr}
\end{itemize}

Πρόσφατες συστηματικές ανασκοπήσεις επιβεβαιώνουν ότι οι \textlatin{DT} εφαρμογές εστιάζουν κυρίως σε αστικές υποδομές και βιομηχανικά περιβάλλοντα [Doulka, 2024], αφήνοντας ανεξερεύνητο τον τομέα των εκπαιδευτικών κτιρίων --- κενό που η παρούσα εργασία καλύπτει.

Τα βασικά ευρήματα της πιλοτικής εφαρμογής συνοψίζονται ως εξής:

\begin{itemize}
\item Ο κλειστός βρόχος ελέγχου λειτουργεί --- αποδείχθηκε αρχιτεκτονικά και επιχειρησιακά
\item Το \textlatin{FridgeGuardian} εντόπισε δύο ανωμαλίες κατανάλωσης (Ιούλιος και Οκτώβριος 2025) και απέστειλε αυτόματη ειδοποίηση --- πρακτικό παράδειγμα Προγνωστικής Συντήρησης
\item Αλλαγή συμπεριφοράς προσωπικού κυλικείου επιβεβαιώθηκε από το ίδιο το προσωπικό και τεκμηριώνεται από τα γραφήματα ανοιγμάτων θύρας
\item Παρατηρήθηκε μείωση 19,7\% στους λογαριασμούς ρεύματος --- στοιχείο που συσχετίζεται με τις παρεμβάσεις του συστήματος χωρίς να αποδεικνύει αιτιακή σχέση, όπως αναλύεται στο Κεφάλαιο 7
\item Επιτυχής υιοθέτηση από μη τεχνικούς χρήστες --- ο διευθυντής του σχολείου ενσωμάτωσε την πλατφόρμα στην καθημερινή του ροή εργασίας χωρίς εξειδικευμένη κατάρτιση
\end{itemize}

Η οικονομική βιωσιμότητα της πλατφόρμας δεν εξαρτάται από την απόδειξη ενεργειακής εξοικονόμησης. Με κόστος υλικών 245,17\euro\ και εκτιμώμενο ετήσιο κόστος συντήρησης $\sim$1.140\euro, η αποφυγή έστω μίας βλάβης ψυκτικού εξοπλισμού ετησίως --- εκτιμώμενο κόστος 500--1.300\euro\ ανά περιστατικό --- καθιστά την επένδυση οικονομικά βιώσιμη. Η κύρια αξία έγκειται στη λειτουργική εποπτεία 24/7 και στην έγκαιρη ανίχνευση αστοχιών --- λειτουργίες που ευθυγραμμίζονται με τους στόχους του \textlatin{DIGITAL EUROPE} για \textlatin{Local Digital Twins} \cite{EuropeanCommission2023} και της Ευρωπαϊκής Πράσινης Συμφωνίας \cite{EUParliament2024}.

Οι κύριοι περιορισμοί --- απουσία υποδιαίρεσης κατανάλωσης ανά υποσύστημα, περιορισμένο αισθητηριακό δίκτυο και εξάρτηση από \textlatin{WiFi} --- αποτελούν οδηγό για μελλοντική έρευνα. Τα επόμενα φυσικά βήματα περιλαμβάνουν:

\begin{itemize}
\item Εγκατάσταση έξυπνων μετρητών ανά υποσύστημα για απόδειξη αιτιακής σχέσης
\item Επέκταση αισθητηριακού δικτύου με αισθητήρες \textlatin{CO2} και πλήρη κάλυψη αιθουσών
\item Ενσωμάτωση \textlatin{BIM} μοντέλων --- ερευνητικά ώριμο πεδίο που έχει αναλυθεί στον ελληνικό ακαδημαϊκό χώρο [Φίλογλου, 2021· Πούλιος, 2023]
\item Εξέλιξη σε \textlatin{Cognitive Digital Twin} με αλγορίθμους ανίχνευσης ανωμαλιών
\end{itemize}

Η πλατφόρμα --- διαθέσιμη στο \textlatin{digitaltwin.pi-tech.gr} --- είναι ήδη έτοιμη να δεχτεί αυτές τις επεκτάσεις, αποτελώντας λειτουργικό σημείο αναφοράς για μελλοντικές υλοποιήσεις \textlatin{Smart Campus} υποδομών.
